\documentclass[aspectratio=169]{beamer}
\usepackage[utf8]{inputenc}
%\usepackage[sfmath]{kpfonts}
\usepackage{helvet}
%\usepackage{lmodern}
\usepackage{siunitx}
\usepackage{tikz}

\usepackage{xcolor}
\usepackage{hyperref}

\hypersetup{
  colorlinks,
  citecolor=cyan,
  linkcolor=blue}


\usetheme{Madrid}
%\usecolortheme{default}

\usepackage[backend=biber]{biblatex}
%\bibliography{refs.bib}
\addbibresource{refs.bib}

%------------------------------------------------------------
%This block of code defines the information to appear in the
%Title page
\title[Microcontroller architecture] %optional
{Microcontroller architecture}

\subtitle{ISA and Microarchitecture}

\author{Daniel Giovanni Martínez Sandoval} 

\institute[CUCEI - UDG]
{
  \textsc{
    \textbf{Universidad de Guadalajara}\\
    {\tiny Centro Universitario de Ciencias Exactas e Ingenierías}
  }
}

\date[September 2024] % (optional)
{\footnotesize I7266 - Programación de Sistemas Embebidos - 2024B }

\titlegraphic { 
  \begin{tikzpicture}[overlay,remember picture]
  \node[right=-2.5cm] at (current page.345){
    \includegraphics[height=20mm]{logo_udg_color.eps}
  };
  \end{tikzpicture}
}
%\logo{\includegraphics[height=15mm]{logo_udg_color.eps}}

%End of title page configuration block
%------------------------------------------------------------



%------------------------------------------------------------
%The next block of commands puts the table of contents at the 
%beginning of each section and highlights the current section:

\AtBeginSection[]
{
  \begin{frame}
    \frametitle{Table of Contents}
    \tableofcontents[currentsection]
\end{frame}
}
%------------------------------------------------------------


\begin{document}

%The next statement creates the title page.
\frame{\titlepage}


%---------------------------------------------------------
%This block of code is for the table of contents after
%the title page
\begin{frame}
\frametitle{Table of Contents}
\tableofcontents
\end{frame}
%---------------------------------------------------------


\section{Computer architecture}
\subsection{What is an architecture?}

%---------------------------------------------------------
\begin{frame}
\frametitle{Definition of architecture}
  \begin{block}{Architecture}
    Architecture describes the internal organization of a computer in an abstract way; that is, it defines the capabilities of the computer and its programming model.\cite{clements} It outlines how a computer's hardware components interact and work together to execute programs and process data.   
  \end{block}
\end{frame}
%---------------------------------------------------------

%---------------------------------------------------------
\begin{frame}
  \frametitle{}
  In embedded systems, architecture plays a crucial role because these systems are mostly designed for specific tasks rather than general-purpose computing. Some things to keep in mind when choosing an architecture are:
  \begin{itemize}
    \item <1-> \textbf{Power consumption:} Many embedded systems operate on battery power. A microcontroller packed with features but with a high current consumption is not always feasible for battery-operated systems. 
    \item <2-> \textbf{Cost and size:} Choosing a microcontroller with just enough features to develop the intended functionality is key to saving money, but also in some cases the size, which in turn also reduces the cost.
    \item <3-> \textbf{Performance and memory:} Depending on the application, some devices might only run certain OS which is targeted for a specific arquitechture, or maybe they need to be compatible with external memory. 
  \end{itemize}
\end{frame}
%---------------------------------------------------------

%---------------------------------------------------------
\begin{frame}
  \frametitle{Definition of architecture}
    \begin{block}{Architecture}
      Architecture describes the internal organization of a computer in an abstract way; that is, it defines the capabilities of the computer and its programming model.\cite{clements} It outlines how a computer's hardware components interact and work together to execute programs and process data.   
    \end{block}
  \end{frame}
  %---------------------------------------------------------

  %---------------------------------------------------------
\begin{frame}
  \frametitle{ISA}
    \begin{block}{Instruction Set Architecture}
      An Instruction Set Architecture (ISA) is part of the abstract model of a computer that defines how the CPU is controlled by the software. The ISA acts as an interface between the hardware and the software, specifying both what the processor is capable of doing as well as how it gets done. \cite{clements}
    \end{block}
    The ISA provides the only way through which a user is able to interact with the hardware. It can be viewed as a programmer's manual because it's the portion of the machine that's visible to the assembly language programmer, the compiler writer, and the application programmer. 
  \end{frame}
  %---------------------------------------------------------

  %---------------------------------------------------------
\begin{frame}
  \frametitle{Definition of architecture}
    \begin{block}{Architecture}
      Architecture describes the internal organization of a computer in an abstract way; that is, it defines the capabilities of the computer and its programming model.\cite{clements} It outlines how a computer's hardware components interact and work together to execute programs and process data.   
    \end{block}
  \end{frame}
  %---------------------------------------------------------

\begin{frame}
    \frametitle{References}
    % This prints the bibliography on the slide
    %\bibliographystyle{IEEEtran} % We choose the "plain" reference style
    %\bibliography{refs} % Entries are in the refs.bib file
    \printbibliography
\end{frame}
%---------------------------------------------------------

\end{document}

% \begin{frame}
%   \frametitle{Format presentation}
  
%   In this slide, some important text will be
%   \alert{highlighted} because it's important.
%   Please, don't abuse it.
  
%   \begin{block}{Remark}
%   Sample text
%   \end{block}
  
%   \begin{alertblock}{Important theorem}
%   Sample text in red box
%   \end{alertblock}
  
%   \begin{examples}
%   Sample text in green box. The title of the block is ``Examples".
%   \end{examples}
% \end{frame}


% %---------------------------------------------------------
% \begin{frame}
%   \frametitle{Frame title}
%   % Body
% \end{frame}
% %---------------------------------------------------------
% %---------------------------------------------------------
% \begin{frame}
%   \frametitle{Embedded systems everywhere}
%   Embedded systems are a little bit hard to define. This is because the applications for embedded systems has changed over time and has become very broad.\\[5mm] \pause
%   If we say that embedded systems are ``small'' computer systems that make part of bigger electronic or mechanical systems, how many devices come to your mind?
% \end{frame}
  
%   %---------------------------------------------------------
%   \begin{frame}
%     \frametitle{Embedded systems everywhere}
%     Some areas full of embedded systems, in fact, are: \pause
%     \begin{columns}
%       \begin{column}{0.5\textwidth}
%         \begin{itemize}
%           \item<2-> Consumer Electronics
%           \item<3-> Home Appliances
%           \item<4-> Automotive
%           \item<5-> Industrial Applications
%           \item<6-> Medical Devices
%           \item<7-> Telecommunications
%         \end{itemize}
%       \end{column}
%       \begin{column}{0.5\textwidth}
%         \begin{center}
%           \begin{itemize}
%             \item<8-> Aerospace and Defense
%             \item<9-> Office Automation
%             \item<10-> Security Systems
%             \item<11-> Smart Home Devices
%             \item<12-> Wearable Devices
%           \end{itemize}
%         \end{center}
%       \end{column}
%     \end{columns}
%   \end{frame}
  
%   %---------------------------------------------------------
  
  
%   %---------------------------------------------------------
%   %Example of the \pause command
%   \begin{frame}
%   In this slide \pause
  
%   the text will be partially visible \pause
  
%   And finally everything will be there
% \end{frame}
%   %---------------------------------------------------------
  
%   \section{Second section}
  
%   %---------------------------------------------------------
%   %Highlighting text
%   \begin{frame}
%   \frametitle{Sample frame title}
  
%   In this slide, some important text will be
%   \alert{highlighted} because it's important.
%   Please, don't abuse it.
  
%   \begin{block}{Remark}
%   Sample text
%   \end{block}
  
%   \begin{alertblock}{Important theorem}
%   Sample text in red box
%   \end{alertblock}
  
%   \begin{examples}
%   Sample text in green box. The title of the block is ``Examples".
%   \end{examples}
% \end{frame}
%   %---------------------------------------------------------
  
  
%   %---------------------------------------------------------
%   %Two columns
%   \begin{frame}
%   \frametitle{Two-column slide}
  
%   \begin{columns}
  
%   \column{0.5\textwidth}
%   This is a text in first column.
%   $$E=mc^2$$
%   \begin{itemize}
%   \item First item
%   \item Second item
%   \end{itemize}
  
%   \column{0.5\textwidth}
%   This text will be in the second column
%   and on a second tought this is a nice looking
%   layout in some cases.
%   \end{columns}
% \end{frame}
%   %---------------------------------------------------------

